% Options for packages loaded elsewhere
\PassOptionsToPackage{unicode}{hyperref}
\PassOptionsToPackage{hyphens}{url}
\documentclass[
]{article}
\usepackage{xcolor}
\usepackage{amsmath,amssymb}
\setcounter{secnumdepth}{-\maxdimen} % remove section numbering
\usepackage{iftex}
\ifPDFTeX
  \usepackage[T1]{fontenc}
  \usepackage[utf8]{inputenc}
  \usepackage{textcomp} % provide euro and other symbols
\else % if luatex or xetex
  \usepackage{unicode-math} % this also loads fontspec
  \defaultfontfeatures{Scale=MatchLowercase}
  \defaultfontfeatures[\rmfamily]{Ligatures=TeX,Scale=1}
\fi
\usepackage{lmodern}
\ifPDFTeX\else
  % xetex/luatex font selection
\fi
% Use upquote if available, for straight quotes in verbatim environments
\IfFileExists{upquote.sty}{\usepackage{upquote}}{}
\IfFileExists{microtype.sty}{% use microtype if available
  \usepackage[]{microtype}
  \UseMicrotypeSet[protrusion]{basicmath} % disable protrusion for tt fonts
}{}
\makeatletter
\@ifundefined{KOMAClassName}{% if non-KOMA class
  \IfFileExists{parskip.sty}{%
    \usepackage{parskip}
  }{% else
    \setlength{\parindent}{0pt}
    \setlength{\parskip}{6pt plus 2pt minus 1pt}}
}{% if KOMA class
  \KOMAoptions{parskip=half}}
\makeatother
\usepackage{color}
\usepackage{fancyvrb}
\newcommand{\VerbBar}{|}
\newcommand{\VERB}{\Verb[commandchars=\\\{\}]}
\DefineVerbatimEnvironment{Highlighting}{Verbatim}{commandchars=\\\{\}}
% Add ',fontsize=\small' for more characters per line
\newenvironment{Shaded}{}{}
\newcommand{\AlertTok}[1]{\textcolor[rgb]{1.00,0.00,0.00}{\textbf{#1}}}
\newcommand{\AnnotationTok}[1]{\textcolor[rgb]{0.38,0.63,0.69}{\textbf{\textit{#1}}}}
\newcommand{\AttributeTok}[1]{\textcolor[rgb]{0.49,0.56,0.16}{#1}}
\newcommand{\BaseNTok}[1]{\textcolor[rgb]{0.25,0.63,0.44}{#1}}
\newcommand{\BuiltInTok}[1]{\textcolor[rgb]{0.00,0.50,0.00}{#1}}
\newcommand{\CharTok}[1]{\textcolor[rgb]{0.25,0.44,0.63}{#1}}
\newcommand{\CommentTok}[1]{\textcolor[rgb]{0.38,0.63,0.69}{\textit{#1}}}
\newcommand{\CommentVarTok}[1]{\textcolor[rgb]{0.38,0.63,0.69}{\textbf{\textit{#1}}}}
\newcommand{\ConstantTok}[1]{\textcolor[rgb]{0.53,0.00,0.00}{#1}}
\newcommand{\ControlFlowTok}[1]{\textcolor[rgb]{0.00,0.44,0.13}{\textbf{#1}}}
\newcommand{\DataTypeTok}[1]{\textcolor[rgb]{0.56,0.13,0.00}{#1}}
\newcommand{\DecValTok}[1]{\textcolor[rgb]{0.25,0.63,0.44}{#1}}
\newcommand{\DocumentationTok}[1]{\textcolor[rgb]{0.73,0.13,0.13}{\textit{#1}}}
\newcommand{\ErrorTok}[1]{\textcolor[rgb]{1.00,0.00,0.00}{\textbf{#1}}}
\newcommand{\ExtensionTok}[1]{#1}
\newcommand{\FloatTok}[1]{\textcolor[rgb]{0.25,0.63,0.44}{#1}}
\newcommand{\FunctionTok}[1]{\textcolor[rgb]{0.02,0.16,0.49}{#1}}
\newcommand{\ImportTok}[1]{\textcolor[rgb]{0.00,0.50,0.00}{\textbf{#1}}}
\newcommand{\InformationTok}[1]{\textcolor[rgb]{0.38,0.63,0.69}{\textbf{\textit{#1}}}}
\newcommand{\KeywordTok}[1]{\textcolor[rgb]{0.00,0.44,0.13}{\textbf{#1}}}
\newcommand{\NormalTok}[1]{#1}
\newcommand{\OperatorTok}[1]{\textcolor[rgb]{0.40,0.40,0.40}{#1}}
\newcommand{\OtherTok}[1]{\textcolor[rgb]{0.00,0.44,0.13}{#1}}
\newcommand{\PreprocessorTok}[1]{\textcolor[rgb]{0.74,0.48,0.00}{#1}}
\newcommand{\RegionMarkerTok}[1]{#1}
\newcommand{\SpecialCharTok}[1]{\textcolor[rgb]{0.25,0.44,0.63}{#1}}
\newcommand{\SpecialStringTok}[1]{\textcolor[rgb]{0.73,0.40,0.53}{#1}}
\newcommand{\StringTok}[1]{\textcolor[rgb]{0.25,0.44,0.63}{#1}}
\newcommand{\VariableTok}[1]{\textcolor[rgb]{0.10,0.09,0.49}{#1}}
\newcommand{\VerbatimStringTok}[1]{\textcolor[rgb]{0.25,0.44,0.63}{#1}}
\newcommand{\WarningTok}[1]{\textcolor[rgb]{0.38,0.63,0.69}{\textbf{\textit{#1}}}}
\usepackage{longtable,booktabs,array}
\newcounter{none} % for unnumbered tables
\usepackage{calc} % for calculating minipage widths
% Correct order of tables after \paragraph or \subparagraph
\usepackage{etoolbox}
\makeatletter
\patchcmd\longtable{\par}{\if@noskipsec\mbox{}\fi\par}{}{}
\makeatother
% Allow footnotes in longtable head/foot
\IfFileExists{footnotehyper.sty}{\usepackage{footnotehyper}}{\usepackage{footnote}}
\makesavenoteenv{longtable}
\setlength{\emergencystretch}{3em} % prevent overfull lines
\providecommand{\tightlist}{%
  \setlength{\itemsep}{0pt}\setlength{\parskip}{0pt}}
\usepackage{bookmark}
\IfFileExists{xurl.sty}{\usepackage{xurl}}{} % add URL line breaks if available
\urlstyle{same}
\hypersetup{
  hidelinks,
  pdfcreator={LaTeX via pandoc}}

\author{}
\date{}

\begin{document}

\section{Learning Whom to Trust with
MACE}\label{learning-whom-to-trust-with-mace}

Dirk Hovy \(^{1}\) Taylor Berg-Kirkpatrick \(^{2}\) Ashish Vaswani
\(^{1}\) Eduard Hovy \(^{3}\)

\begin{enumerate}
\def\labelenumi{(\arabic{enumi})}
\tightlist
\item
  Information Sciences Institute, University of Southern California,
  Marina del Rey\\
\item
  Computer Science Division, University of California at Berkeley\\
\item
  Language Technology Institute, Carnegie Mellon University, Pittsburgh
  \{dirkh, avaswani\}@isi.edu, tberg@cs.berkeley.edu, hovy@cmu.edu
\end{enumerate}

\section{Abstract}\label{abstract}

Non-expert annotation services like Amazon's Mechanical Turk (AMT) are
cheap and fast ways to evaluate systems and provide categorical
annotations for training data. Unfortunately, some annotators choose bad
labels in order to maximize their pay. Manual identification is tedious,
so we experiment with an item-response model. It learns in an
unsupervised fashion to a) identify which annotators are trustworthy and
b) predict the correct underlying labels. We match performance of more
complex state-of-the-art systems and perform well even under adversarial
conditions. We show considerable improvements over standard baselines,
both for predicted label accuracy and trustworthiness estimates. The
latter can be further improved by introducing a prior on model
parameters and using Variational Bayes inference. Additionally, we can
achieve even higher accuracy by focusing on the instances our model is
most confident in (trading in some recall), and by incorporating
annotated control instances. Our system, MACE (Multi-Annotator
Competence Estimation), is available for download \(^{1}\) .

\section{1 Introduction}\label{introduction}

Amazon's MechanicalTurk (AMT) is frequently used to evaluate experiments
and annotate data in NLP (Callison-Burch et al., 2010; Callison-Burch
and Dredze, 2010; Jha et al., 2010; Zaidan and Callison-Burch, 2011).
However, some turkers try to maximize their pay by supplying quick
answers that have nothing to do with the correct label. We refer to this
type of annotator as a spammer. In order to mitigate the effect of
spammers, researchers typically collect multiple annotations of the same
instance so that they can, later, use de-noising methods to infer the
best label. The simplest approach is majority voting, which weights all
answers equally. Unfortunately, it is easy for majority voting to go
wrong. A common and simple spammer strategy for categorical labeling
tasks is to always choose the same (often the first) label. When
multiple spammers follow this strategy, the majority can be incorrect.
While this specific scenario might seem simple to correct for (remove
annotators that always produce the same label), the situation grows more
tricky when spammers do not annotate consistently, but instead choose
labels at random. A more sophisticated approach than simple majority
voting is required.

If we knew whom to trust, and when, we could reconstruct the correct
labels. Yet, the only way to be sure we know whom to trust is if we knew
the correct labels ahead of time. To address this circular problem, we
build a generative model of the annotation process that treats the
correct labels as latent variables. We then use unsupervised learning to
estimate parameters directly from redundant annotations. This is a
common approach in the class of unsupervised models called item-response
models (Dawid and Skene, 1979; Whitehill et al., 2009; Carpenter, 2008;
Raykar and Yu, 2012). While such models have been implemented in other
fields (e.g., vision), we are not aware of their availability for NLP
tasks (see also Section 6).

Our model includes a binary latent variable that explicitly encodes if
and when each annotator is spamming, as well as parameters that model
the annotator's specific spamming ``strategy''. Impor-

tantly, the model assumes that labels produced by an annotator when
spamming are independent of the true label (though, a spammer can still
produce the correct label by chance).

In experiments, our model effectively differentiates dutiful annotators
from spammers (Section 4), and is able to reconstruct the correct label
with high accuracy (Section 5), even under extremely adversarial
conditions (Section 5.2). It does not require any annotated instances,
but is capable of including varying levels of supervision via token
constraints (Section 5.2). We consistently outperform majority voting,
and achieve performance equal to that of more complex state-of-the-art
models. Additionally, we find that thresholding based on the posterior
label entropy can be used to trade off coverage for accuracy in label
reconstruction, giving considerable gains (Section 5.1). In tasks where
correct answers are more important than answering every instance, e.g.,
when constructing a new annotated corpus, this feature is extremely
valuable. Our contributions are:

\begin{itemize}
\tightlist
\item
  We demonstrate the effectiveness of our model on real world AMT
  datasets, matching the accuracy of more complex state-of-the-art
  systems\\
\item
  We show how posterior entropy can be used to trade some coverage for
  considerable gains in accuracy\\
\item
  We study how various factors affect performance, including number of
  annotators, annotator strategy, and available supervision\\
\item
  We provide MACE (Multi-Annotator Competence Estimation), a Java-based
  implementation of a simple and scalable unsupervised model that
  identifies malicious annotators and predicts labels with high accuracy
\end{itemize}

\section{2 Model}\label{model}

We keep our model as simple as possible so that it can be effectively
trained from data where annotator quality is unknown. If the model has
too many parameters, unsupervised learning can easily pick up on and
exploit coincidental correlations in the data. Thus, we make a modeling
assumption that keeps our parameterization simple. We assume that an
annotator always produces the correct label when

\textit{[Image omitted: images/0f77943c1131a0feb5adf1ef97de95191dce8288c38149cc3c4851d19e9ac6c3.jpg]}

flowchart

\begin{Shaded}
\begin{Highlighting}[]
\NormalTok{graph TD}
\NormalTok{    Ti["Ti"] {-}{-}\textgreater{} Aj["Aij"]}
\NormalTok{    Sij["Sij"] {-}{-}\textgreater{} Aj}
\NormalTok{    style Ti fill:\#fff,stroke:\#000}
\NormalTok{    style Sij fill:\#fff,stroke:\#000}
\NormalTok{    style Aj fill:\#000,stroke:\#000}
\NormalTok{    style M fill:\#fff,stroke:\#000}
\NormalTok{    style N fill:\#fff,stroke:\#000}
\end{Highlighting}
\end{Shaded}

Figure 1: Graphical model: Annotator j produces label \(A_{ij}\) on
instance i. Label choice depends on instance's true label \(T_{i}\) ,
and whether j is spamming on i, modeled by binary variable \(S_{ij}\) .
\(N = |instances|\) , \(M = |annotators|\) .

\[
\text { for } i = 1 \dots N:
\]

\[
T _ {i} \sim \text { Uniform }
\]

\[
\text { for } j = 1 \dots M:
\]

\[
S _ {i j} \sim \text { Bernoulli } (1 - \theta_ {j})
\]

\[
\text { if } S _ {i j} = 0:
\]

\[
A _ {i j} = T _ {i}
\]

\[
\mathrm{else}:
\]

\[
A _ {i j} \sim \operatorname{Multinomial} (\xi_ {j})
\]

Figure 2: Generative process: see text for description.

he tries to. While this assumption does not reflect the reality of AMT,
it allows us to focus the model's power where it's important: explaining
away labels that are not correlated with the correct label.

Our model generates the observed annotations as follows: First, for each
instance i, we sample the true label \(T_{i}\) from a uniform prior.
Then, for each annotator j we draw a binary variable \(S_{ij}\) from a
Bernoulli distribution with parameter \(1 - \theta_{j}\) . \(S_{ij}\)
represents whether or not annotator j is spamming on instance i. We
assume that when an annotator is not spamming on an instance,
i.e.~\(S_{ij} = 0\) , he just copies the true label to produce
annotation \(A_{ij}\) . If \(S_{ij} = 1\) , we say that the annotator is
spamming on the current instance, and \(A_{ij}\) is sampled from a
multinomial with parameter vector \(\xi_{j}\) . Note that in this case
the annotation \(A_{ij}\) does not depend on the true label \(T_{i}\) .
The annotations \(A_{ij}\) are observed,

while the true labels \(T_{i}\) and the spamming indicators \(S_{ij}\)
are unobserved. The graphical model is shown in Figure 1 and the
generative process is described in Figure 2.

The model parameters are \(\theta_{j}\) and \(\xi_{j}\) . \(\theta_{j}\)
specifies the probability of trustworthiness for annotator j (i.e.~the
probability that he is not spamming on any given instance). The learned
value of \(\theta_{j}\) will prove useful later when we try to identify
reliable annotators (see Section 4). The vector \(\xi_{j}\) determines
how annotator j behaves when he is spamming. An annotator can produce
the correct answer even while spamming, but this can happen only by
chance since the annotator must use the same multinomial parameters
\(\xi_{j}\) across all instances. This means that we only learn
annotator biases that are not correlated with the correct label, e.g.,
the strategy of the spammer who always chooses a certain label. This
contrasts with previous work where additional parameters are used to
model the biases that even dutiful annotators exhibit. Note that an
annotator can also choose not to answer, which we can naturally
accommodate because the model is generative. We enhance our generative
model by adding Beta and Dirichlet priors on \(\theta_{j}\) and
\(\xi_{j}\) respectively which allows us to incorporate prior beliefs
about our annotators (section 2.1).

\section{2.1 Learning}\label{learning}

We would like to set our model parameters to maximize the probability of
the observed data, i.e., the marginal data likelihood:

\[
\begin{array}{l} P (A; \theta , \xi) = \\ \sum_ {T, S} \left[ \prod_ {i = 1} ^ {N} P (T _ {i}) \cdot \prod_ {j = 1} ^ {M} P (S _ {i j}; \theta_ {j}) \cdot P (A _ {i j} | S _ {i j}, T _ {i}; \xi_ {j}) \right] \\ \end{array}
\]

where A is the matrix of annotations, S is the matrix of competence
indicators, and T is the vector of true labels.

We maximize the marginal data likelihood using Expectation Maximization
(EM) (Dempster et al., 1977), which has successfully been applied to
similar problems (Dawid and Skene, 1979). We initialize EM randomly and
run for 50 iterations. We perform 100 random restarts, and keep the
model with the best marginal data likelihood. We smooth the M-step by
adding a fixed value \(\delta\) to the fractional counts before
normalizing (Eisner, 2002). We find that smoothing improves accuracy,
but, overall, learning is robust to varying \(\delta\) , and set
\(\delta = \frac{0.1}{num\ labels}\) .

We observe, however, that the average annotator proficiency is usually
high, i.e., most annotators answer correctly. The distribution learned
by EM, however, is fairly linear. To improve the correlation between
model estimates and true annotator proficiency, we would like to add
priors about the annotator behavior into the model. A straightforward
approach is to employ Bayesian inference with Beta priors on the
proficiency parameters, \(\theta_{j}\) . We thus also implement
Variational-Bayes (VB) training with symmetric Beta priors on
\(\theta_{j}\) and symmetric Dirichlet priors on the strategy
parameters, \(\xi_{j}\) . Setting the shape parameters of the Beta
distribution to 0.5 favors the extremes of the distribution, i.e.,
either an annotator tried to get the right answer, or simply did not
care, but (almost) nobody tried ``a little''. With VB training, we
observe improved correlations over all test sets with no loss in
accuracy. The hyper-parameters of the Dirichlet distribution on
\(\xi_{j}\) were clamped to 10.0 for all our experiments with VB
training. Our implementation is similar to Johnson (2007), which the
reader can refer to for details.

\section{3 Experiments}\label{experiments}

We evaluate our method on existing annotated datasets from various AMT
tasks. However, we also want to ensure that our model can handle
adversarial conditions. Since we have no control over the factors in
existing datasets, we create synthetic data for this purpose.

\section{3.1 Natural Data}\label{natural-data}

In order to evaluate our model, we use the datasets from (Snow et al.,
2008) that use discrete label values (some tasks used continuous values,
which we currently do not model). Since they compared AMT annotations to
experts, gold annotations exist for these sets. We can thus evaluate the
accuracy of the model as well as the proficiency of each annotator. We
show results for word sense disambiguation (WSD: 177 items, 34
annotators), recognizing textual entailment (RTE: 800 items, 164
annotators), and recognizing temporal relation (Temporal: 462 items, 76
annotators).

\section{3.2 Synthetic Data}\label{synthetic-data}

In addition to the datasets above, we generate synthetic data in order
to control for different

factors. This also allows us to create a gold standard to which we can
compare. We generate data sets with 100 items, using two or four
possible labels.

For each item, we generate answers from 20 different annotators. The
``annotators'' are functions that return one of the available labels
according to some strategy. Better annotators have a smaller chance of
guessing at random.

For various reasons, usually not all annotators see or answer all items.
We thus remove a randomly selected subset of answers such that each item
is only answered by 10 of the annotators. See Figure 3 for an example
annotation of three items.

annotators

items

-

0

0

1

-

0

-

-

0

-

1

-

-

0

-

1

0

-

-

0

-

-

0

-

0

1

-

0

-

0

Figure 3: Annotations: 10 annotators on three items, labels \(\{1, 0\}\)
, 5 annotations/item. Missing annotations marked `--'

\section{3.3 Evaluations}\label{evaluations}

First, we want to know which annotators to trust. We evaluate whether
our model's learned trustworthiness parameters \(\theta_{j}\) can be
used to identify these individuals (Section 4).

We then compare the label predicted by our model and by majority voting
to the correct label. The results are reported as accuracy (Section 5).
Since our model computes posterior entropies for each instance, we can
use this as an approximation for the model's confidence in the
prediction. If we focus on predictions with high confidence (i.e., low
entropy), we hope to see better accuracy, even at the price of leaving
some items unanswered. We evaluate this trade-off in Section 5.1. In
addition, we investigate the influence of the number of spammers and
their strategy on the accuracy of our model (Section 5.2).

\section{4 Identifying Reliable
Annotators}\label{identifying-reliable-annotators}

One of the distinguishing features of the model is that it uses a
parameter for each annotator to estimate whether or not they are
spamming. Can we use this parameter to identify trustworthy individuals,
to invite them for future tasks, and block untrustworthy ones?

RTE

Temporal

WSD

raw agreement

0.78

0.73

0.81

Cohen\textquotesingle s κ

0.70

0.80

0.13

G-index

0.76

0.73

0.81

MACE-EM

0.87

0.88

0.44

MACE-VB (0.5,0.5)

0.91

0.90

0.90

Table 1: Correlation with annotator proficiency: Pearson \(\rho\) of
different methods for various data sets. MACE-VB's trustworthiness
parameter (trained with Variational Bayes with \(\alpha = \beta = 0.5\)
) correlates best with true annotator proficiency.

It is natural to apply some form of weighting. One approach is to assume
that reliable annotators agree more with others than random annotators.
Inter-annotator agreement is thus a good candidate to weigh the answers.
There are various measures for inter-annotator agreement.

Tratz and Hovy (2010) compute the average agreement of each annotator
and use it as a weight to identify reliable ones. Raw agreement can be
directly computed from the data. It is related to majority voting, since
it will produce high scores for all members of the majority class. Raw
agreement is thus a very simple measure.

In contrast, Cohen's \(\kappa\) corrects the agreement between two
annotators for chance agreement. It is widely used for inter-annotator
agreement in annotation tasks. We also compute the \(\kappa\) values for
each pair of annotators, and average them for each annotator (similar to
the approach in Tratz and Hovy (2010)). However, whenever one label is
more prevalent (a common case in NLP tasks), \(\kappa\) overestimates
the effect of chance agreement (Feinstein and Cicchetti, 1990) and
penalizes disproportionately. The G-index (Gwet, 2008) corrects for the
number of labels rather than chance agreement.

We compare these measures to our learned trustworthiness parameters
\(\theta_{j}\) in terms of their ability to select reliable annotators.
A better measure should lend higher score to annotators who answer
correctly more often than others. We thus compare the ratings of each
measure to the true proficiency of each annotator. This is the
percentage of annotated items the annotator answered correctly. Methods
that can identify reliable annotators should highly correlate

to the annotator's proficiency. Since the methods use different scales,
we compute Pearson's \(\rho\) for the correlation coefficient, which is
scale-invariant. The correlation results are shown in Table 1.

The model's \(\theta_{j}\) correlates much more strongly with annotator
proficiency than either \(\kappa\) or raw agreement. The variant trained
with VB performs consistently better than standard EM training, and
yields the best results. This show that our model detects reliable
annotators much better than any of the other measures, which are only
loosely correlated to annotator proficiency.

The numbers for WSD also illustrate the low \(\kappa\) score resulting
when all annotators (correctly) agree on a small number of labels.
However, all inter-annotator agreement measures suffer from an even more
fundamental problem: removing/ignoring annotators with low agreement
will always improve the overall score, irrespective of the quality of
their annotations. Worse, there is no natural stopping point: deleting
the most egregious outlier always improves agreement, until we have only
one annotator with perfect agreement left (Hovy, 2010). In contrast,
MACE does not discard any annotators, but weighs their contributions
differently. We are thus not losing information. This works well even
under adversarial conditions (see Section 5.2).

5 Recovering the Correct Answer

RTE

Temporal

WSD

majority

0.90

0.93

0.99

Raykar/Yu 2012

0.93

0.94

---

Carpenter 2008

0.93

---

---

MACE-EM/VB

0.93

0.94

0.99

MACE-EM@90

0.95

0.97

0.99

MACE-EM@75

0.95

0.97

1.0

MACE-VB@90

0.96

0.97

1.0

MACE-VB@75

0.98

0.98

1.0

Table 2: Accuracy of different methods on data sets from (Snow et al.,
2008). MACE-VB uses Variational Bayes training. Results @n use the n\%
items the model is most confident in (Section 5.1). Results below double
line trade coverage for accuracy and are thus not comparable to upper
half.

The previous sections showed that our model reliably identifies
trustworthy annotators. However, we also want to find the most likely
correct answer. Using majority voting often fails to find the correct
label. This problem worsens when there are more than two labels. We need
to take relative majorities into account or break ties when two or more
labels receive the same number of votes. This is deeply unsatisfying.

Figure 2 shows the accuracy of our model on various data sets from Snow
et al.~(2008). The model outperforms majority voting on both RTE and
Temporal recognition sets. It performs as well as majority voting for
the WSD task. This last set is somewhat of an exception, though, since
almost all annotators are correct all the time, so majority voting is
trivially correct. Still, we need to ensure that the model does not
perform worse under such conditions. The results for RTE and Temporal
data also rival those reported in Raykar and Yu (2012) and Carpenter
(2008), yet were achieved with a much simpler model.

Carpenter (2008) models instance difficulty as a parameter. While it
seems intuitively useful to model which items are harder than other, it
increases the parameter space more than our trustworthiness variable. We
achieve comparable performance without modeling difficulty, which
greatly simplifies inference. The model of Raykar and Yu (2012) is more
similar to our approach, in that it does not model item difficulty.
However, it adds an extra step that learns priors from the estimated
parameters. In our model, this is part of the training process. For more
details on both models, see Section 6.

\section{5.1 Trading Coverage for
Accuracy}\label{trading-coverage-for-accuracy}

Sometimes, we want to produce an answer for every item (e.g., when
evaluating a data set), and sometimes, we value good answers more than
answering all items (e.g., when developing an annotated corpus). Jha et
al.~(2010) have demonstrated how to achieve better coverage (i.e.,
answer more items) by relaxing the majority voting constraints.
Similarly, we can improve accuracy if we only select high quality
annotations, even if this incurs lower coverage.

We provide a parameter in MACE that allows users to set a threshold for
this trade-off: the model only returns a label for an instance if it is
sufficiently confident in its answer. We approximate the model's
confidence by the posterior entropy of

\textit{[Image omitted: images/395d48af4902a6ed6d678832a5d25742c78f890a34d93a53462fbc8054da732e.jpg]}

line

{\def\LTcaptype{none} % do not increment counter
\begin{longtable}[]{@{}llll@{}}
\toprule\noalign{}
coverage & MACE-EM & MACE-VB & majority \\
\midrule\noalign{}
\endhead
\bottomrule\noalign{}
\endlastfoot
0.5 & 0.985 & 0.985 & 0.983 \\
0.6 & 0.973 & 0.986 & 0.975 \\
0.7 & 0.973 & 0.982 & 0.963 \\
0.8 & 0.968 & 0.975 & 0.955 \\
0.9 & 0.955 & 0.960 & 0.933 \\
1.0 & 0.930 & 0.927 & 0.898 \\
\end{longtable}
}

\textit{[Image omitted: images/27e0f72364720a2e6f57ad7e9eec2f63dc776d2af72f49844090e8654ae6733c.jpg]}

line

{\def\LTcaptype{none} % do not increment counter
\begin{longtable}[]{@{}llll@{}}
\toprule\noalign{}
coverage & Series 1 & Series 2 & Series 3 \\
\midrule\noalign{}
\endhead
\bottomrule\noalign{}
\endlastfoot
0.5 & 0.97 & 0.97 & 0.97 \\
0.6 & 0.97 & 0.97 & 0.98 \\
0.7 & 0.97 & 0.97 & 0.98 \\
0.8 & 0.97 & 0.97 & 0.98 \\
0.9 & 0.96 & 0.96 & 0.97 \\
1.0 & 0.94 & 0.94 & 0.93 \\
\end{longtable}
}

Figure 4: Tradeoff between coverage and accuracy for RTE (left) and
temporal (right). Lower thresholds lead to less coverage, but result in
higher accuracy.

each instance. However, entropy depends strongly on the specific makeup
of the dataset (number of annotators and labels, etc.), so it is hard
for the user to set a specific threshold.

Instead of requiring an exact entropy value, we provide a simple
thresholding between 0.0 and 1.0 (setting the threshold to 1.0 will
include all items). After training, MACE orders the posterior entropies
for all instances and selects the value that covers the selected
fraction of the instances. The threshold thus roughly corresponds to
coverage. It then only returns answers for instances whose entropy is
below the threshold. This procedure is similar to precision/recall
curves.

Jha et al.~(2010) showed the effect of varying the relative majority
required, i.e., requiring that at least n out of 10 annotators have to
agree to count an item. We use that method as baseline comparison,
evaluating the effect on coverage and accuracy when we vary n from 5 to
10.

Figure 4 shows the tradeoff between coverage and accuracy for two data
sets. Lower thresholds produce more accurate answers, but result in
lower coverage, as some items are left blank. If we produce answers for
all items, we achieve accuracies of 0.93 for RTE and 0.94 for Temporal,
but by excluding just the 10\% of items in which the model is least
confident, we achieve accuracies as high as 0.95 for RTE and 0.97 for
Temporal. We omit the results for WSD here, since there is little
headroom and they are thus not very informative. Using Variational Bayes
inference consistently achieves higher results for the same coverage
than the standard implementation. Increasing the required majority also
improves accuracy, although not as much, and the loss in coverage is
larger and cannot be controlled. In contrast, our method allows us to
achieve better accuracy at a smaller, controlled loss in coverage.

\section{5.2 Influence of Strategy, Number of Annotators, and
Supervision}\label{influence-of-strategy-number-of-annotators-and-supervision}

Adverse Strategy We showed that our model recovers the correct answer
with high accuracy. However, to test whether this is just a function of
the annotator pool, we experiment with varying the trustworthiness of
the pool. If most annotators answer correctly, majority voting is
trivially correct, as is our model. What happens, however, if more and
more annotators are unreliable? While some agreement can arise from
randomness, majority voting is bound to become worse---can our model
overcome this problem? We set up a second set of experiments to test
this, using synthetic data. We choose 20 annotators and vary the amount
of good annotators among them from 0 to 10 (after which the trivial case
sets in). We define a good annotator as one who answers correctly 95\%
of the time. \(^{2}\) Adverse annotators select their answers randomly
or always choose a certain value (minimal annotators). These are two
frequent strategies of spammers.

For different numbers of labels and varying percentage of spammers, we
measure the accuracy

\textit{[Image omitted: images/367853605ba19f1db0b142eca83ec59ecfd692ca30670c609c5d519442abe3fe.jpg]}

line

{\def\LTcaptype{none} % do not increment counter
\begin{longtable}[]{@{}
  >{\raggedright\arraybackslash}p{(\linewidth - 6\tabcolsep) * \real{0.2824}}
  >{\raggedright\arraybackslash}p{(\linewidth - 6\tabcolsep) * \real{0.2471}}
  >{\raggedright\arraybackslash}p{(\linewidth - 6\tabcolsep) * \real{0.2235}}
  >{\raggedright\arraybackslash}p{(\linewidth - 6\tabcolsep) * \real{0.2471}}@{}}
\toprule\noalign{}
\begin{minipage}[b]{\linewidth}\raggedright
\% of adverse annotators
\end{minipage} & \begin{minipage}[b]{\linewidth}\raggedright
accuracy (blue line)
\end{minipage} & \begin{minipage}[b]{\linewidth}\raggedright
accuracy (red line)
\end{minipage} & \begin{minipage}[b]{\linewidth}\raggedright
accuracy (green line)
\end{minipage} \\
\midrule\noalign{}
\endhead
\bottomrule\noalign{}
\endlastfoot
100 & 0.5 & 0.5 & 0.5 \\
90 & 0.5 & 0.5 & 0.5 \\
80 & 0.5 & 0.5 & 0.5 \\
70 & 1.0 & 0.5 & 0.5 \\
60 & 1.0 & 0.6 & 0.5 \\
50 & 1.0 & 0.7 & 1.0 \\
\end{longtable}
}

\textit{[Image omitted: images/ae8ad600a9d79a28d29872a0a07801c20dbec102a37d2dd595e31bde3de74a96.jpg]}

line

{\def\LTcaptype{none} % do not increment counter
\begin{longtable}[]{@{}
  >{\raggedright\arraybackslash}p{(\linewidth - 6\tabcolsep) * \real{0.2857}}
  >{\raggedright\arraybackslash}p{(\linewidth - 6\tabcolsep) * \real{0.2381}}
  >{\raggedright\arraybackslash}p{(\linewidth - 6\tabcolsep) * \real{0.2262}}
  >{\raggedright\arraybackslash}p{(\linewidth - 6\tabcolsep) * \real{0.2500}}@{}}
\toprule\noalign{}
\begin{minipage}[b]{\linewidth}\raggedright
\% of adverse annotators
\end{minipage} & \begin{minipage}[b]{\linewidth}\raggedright
accuracy (blue line)
\end{minipage} & \begin{minipage}[b]{\linewidth}\raggedright
accuracy (red line)
\end{minipage} & \begin{minipage}[b]{\linewidth}\raggedright
accuracy (green line)
\end{minipage} \\
\midrule\noalign{}
\endhead
\bottomrule\noalign{}
\endlastfoot
100 & 0.5 & 0.5 & 0.5 \\
90 & 0.5 & 0.5 & 0.5 \\
80 & 0.5 & 0.5 & 0.5 \\
70 & 1.0 & 0.5 & 0.5 \\
60 & 1.0 & 0.6 & 0.5 \\
50 & 1.0 & 0.7 & 1.0 \\
\end{longtable}
}

\textit{[Image omitted: images/e82bf3f28d4d029d30fd1fd7e59c89b38d152c29613a6f3d009a4f10ab63de50.jpg]}

line

{\def\LTcaptype{none} % do not increment counter
\begin{longtable}[]{@{}llll@{}}
\toprule\noalign{}
\% of adverse annotators & MACE-EM & MACE-VB & majority \\
\midrule\noalign{}
\endhead
\bottomrule\noalign{}
\endlastfoot
100 & 0.25 & 0.25 & 0.25 \\
90 & 0.80 & 0.45 & 0.35 \\
80 & 0.95 & 0.90 & 0.45 \\
70 & 0.98 & 0.98 & 0.68 \\
60 & 0.99 & 0.99 & 0.82 \\
50 & 0.99 & 0.99 & 0.90 \\
\end{longtable}
}

\textit{[Image omitted: images/da02f02871485ea2f83a4230fc4f772004d9bffa214dcc4e844d5941fd1fd6e4.jpg]}

line

{\def\LTcaptype{none} % do not increment counter
\begin{longtable}[]{@{}
  >{\raggedright\arraybackslash}p{(\linewidth - 6\tabcolsep) * \real{0.3478}}
  >{\raggedright\arraybackslash}p{(\linewidth - 6\tabcolsep) * \real{0.2174}}
  >{\raggedright\arraybackslash}p{(\linewidth - 6\tabcolsep) * \real{0.2029}}
  >{\raggedright\arraybackslash}p{(\linewidth - 6\tabcolsep) * \real{0.2319}}@{}}
\toprule\noalign{}
\begin{minipage}[b]{\linewidth}\raggedright
\% of adverse annotators
\end{minipage} & \begin{minipage}[b]{\linewidth}\raggedright
accuracy (blue)
\end{minipage} & \begin{minipage}[b]{\linewidth}\raggedright
accuracy (red)
\end{minipage} & \begin{minipage}[b]{\linewidth}\raggedright
accuracy (green)
\end{minipage} \\
\midrule\noalign{}
\endhead
\bottomrule\noalign{}
\endlastfoot
100 & 0.25 & 0.25 & 0.25 \\
90 & 0.25 & 0.25 & 0.25 \\
80 & 0.25 & 0.25 & 0.25 \\
70 & 0.40 & 0.25 & 0.25 \\
60 & 0.98 & 0.40 & 0.25 \\
50 & 1.00 & 0.60 & 1.00 \\
\end{longtable}
}

Figure 5: Influence of adverse annotator strategy on label accuracy
(y-axis). Number of possible labels varied between 2 (top row) and 4
(bottom row). Adverse annotators either choose at random (a) or always
select the first label (b). MACE needs fewer good annotators to recover
the correct answer.

\textit{[Image omitted: images/ced5be4f8568581d8f0a9a3bf8988fd4212357071498ae55a43961828bebe1dd.jpg]}

line

{\def\LTcaptype{none} % do not increment counter
\begin{longtable}[]{@{}ll@{}}
\toprule\noalign{}
RTE & accuracy \\
\midrule\noalign{}
\endhead
\bottomrule\noalign{}
\endlastfoot
1 & 0.81 \\
2 & 0.77 \\
3 & 0.85 \\
4 & 0.88 \\
5 & 0.90 \\
6 & 0.91 \\
7 & 0.92 \\
8 & 0.92 \\
9 & 0.93 \\
10 & 0.93 \\
\end{longtable}
}

\textit{[Image omitted: images/ae5188411c48a9b591bd0ee73d13925a7a8dd6c8f6e9685891e060230db35fdd.jpg]}

line

{\def\LTcaptype{none} % do not increment counter
\begin{longtable}[]{@{}llll@{}}
\toprule\noalign{}
number of annotators & Line 1 & Line 2 & Line 3 \\
\midrule\noalign{}
\endhead
\bottomrule\noalign{}
\endlastfoot
1 & 0.75 & 0.75 & 0.75 \\
2 & 0.83 & 0.85 & 0.74 \\
3 & 0.89 & 0.90 & 0.83 \\
4 & 0.91 & 0.92 & 0.83 \\
5 & 0.93 & 0.93 & 0.89 \\
6 & 0.93 & 0.93 & 0.89 \\
7 & 0.93 & 0.93 & 0.92 \\
8 & 0.94 & 0.94 & 0.93 \\
9 & 0.94 & 0.94 & 0.94 \\
10 & 0.94 & 0.94 & 0.94 \\
\end{longtable}
}

\textit{[Image omitted: images/cd159d9ae8d389b75133da90d2ee9064b497bbe306726b82c58515111cc72229.jpg]}

line

{\def\LTcaptype{none} % do not increment counter
\begin{longtable}[]{@{}llll@{}}
\toprule\noalign{}
WSD & MACE-EM & MACE-VB & majority \\
\midrule\noalign{}
\endhead
\bottomrule\noalign{}
\endlastfoot
1 & 0.985 & 0.985 & 0.985 \\
2 & 0.990 & 0.993 & 0.985 \\
3 & 0.990 & 0.994 & 0.994 \\
4 & 0.991 & 0.992 & 0.993 \\
5 & 0.993 & 0.994 & 0.994 \\
6 & 0.991 & 0.994 & 0.994 \\
7 & 0.993 & 0.994 & 0.994 \\
8 & 0.994 & 0.994 & 0.994 \\
9 & 0.994 & 0.994 & 0.994 \\
10 & 0.994 & 0.994 & 0.994 \\
\end{longtable}
}

Figure 6: Varying number of annotators: effect on prediction accuracy.
Each point averaged over 10 runs. Note different scale for WSD.

of our model and majority voting on 100 items, averaged over 10 runs for
each condition. Figure 5 shows the effect of annotator proficiency on
both majority voting and our method for both kinds of spammers.
Annotator pool strategy affects majority voting more than our model.
Even with few good annotators, our model learns to dismiss the spammers
as noise. There is a noticeable point on each graph where MACE diverges
from the majority voting line. It thus reaches good accuracy much

faster than majority voting, i.e., with fewer good annotators. This
divergence point happens earlier with more label values when adverse
annotators label randomly. In general, random annotators are easier to
deal with than the ones always choosing the first label. Note that in
cases where we have a majority of adversarial annotators, VB performs
worse than EM, since this condition violates the implicit assumptions we
encoded with the priors in VB. Under these conditions, setting different
priors to reflect the annotator pool should improve performance.

Obviously, both of these pools are extremes: it is unlikely to have so
few good or so many malicious annotators. Most pools will be somewhere
in between. It does show, however, that our model can pick up on
reliable annotators even under very unfavorable conditions. The result
has a practical upshot: AMT allows us to require a minimum rating for
annotators to work on a task. Higher ratings improve annotation quality,
but delay completion, since there are fewer annotators with high
ratings. The results in this section suggest that we can find the
correct answer even in annotator pools with low overall proficiency. We
can thus waive the rating requirement and allow more annotators to work
on the task. This considerably speeds up completion.

Number of Annotators Figure 6 shows the effect different numbers of
annotators have on accuracy. As we increase the number of annotators,
MACE and majority voting achieve better accuracy results. We note that
majority voting results level or even drop when going from an odd to an
even number. In these cases, the new annotator does not improve accuracy
if it goes with the previous majority (i.e., going from 3:2 to 4:2), but
can force an error when going against the previous majority (i.e., from
3:2 to 3:3), by creating a tie. MACE-EM and MACE-VB dominate majority
voting for RTE and Temporal. For WSD, the picture is less clear, where
majority voting dominates when there are fewer annotators. Note that the
differences are minute, though (within 1 percentage point). For very
small pool sizes (\textless{} 3), MACE-VB outperforms both other
methods.

Amount of Supervision So far, we have treated the task as completely
unsupervised. MACE does not require any expert annotations in order to
achieve high accuracy. However, we often have annotations for some of
the items. These annotated data points are usually used as control items
(by removing annotators that answer them incorrectly). If such annotated
data is available, we would like to make use of it. We include an option
that lets users supply annotations for some of the items, and use this
information as token constraints in the E-step of training. In those
cases, the model does not need to estimate the correct value, but only
has to adjust the trust parameter. This leads to improved performance.
\(^{3}\)

We explore for RTE and Temporal how performance changes when we vary the
amount of supervision in increments of 5\%. \(^{4}\) We average over 10
runs for each value of n, each time supplying annotations for a random
set of n items. The baseline uses the annotated label whenever supplied,
otherwise the majority vote, with ties split at random.

Figure 7 shows that, unsurprisingly, all methods improve with additional
supervision, ultimately reaching perfect accuracy. However, MACE uses
the information more effectively, resulting in higher accuracy for a
given amount of supervision. This gain is more pronounced when only
little supervision is available.

\section{6 Related Research}\label{related-research}

Snow et al.~(2008) and Sorokin and Forsyth (2008) showed that Amazon's
MechanicalTurk use in providing non-expert annotations for NLP tasks.
Various models have been proposed for predicting correct annotations
from noisy non-expert annotations and for estimating annotator
trustworthiness. These models divide naturally into two categories:
those that use expert annotations for supervised learning (Snow et al.,
2008; Bian et al., 2009), and completely unsupervised ones. Our method
falls into the latter category because it learns from the redundant
non-expert annotations themselves, and makes no use of expertly
annotated data.

Most previous work on unsupervised models belongs to a class called
``Item-response models'', used in psychometrics. The approaches differ
with respect to which aspect of the annotation process

\textit{[Image omitted: images/0d1bcaaaefd73f96e20ce7b3920963db6e4fd929c82e5252ffdaf99e85f9c8d7.jpg]}

line

{\def\LTcaptype{none} % do not increment counter
\begin{longtable}[]{@{}
  >{\raggedright\arraybackslash}p{(\linewidth - 6\tabcolsep) * \real{0.3182}}
  >{\raggedright\arraybackslash}p{(\linewidth - 6\tabcolsep) * \real{0.2121}}
  >{\raggedright\arraybackslash}p{(\linewidth - 6\tabcolsep) * \real{0.2273}}
  >{\raggedright\arraybackslash}p{(\linewidth - 6\tabcolsep) * \real{0.2424}}@{}}
\toprule\noalign{}
\begin{minipage}[b]{\linewidth}\raggedright
\% of supervised data
\end{minipage} & \begin{minipage}[b]{\linewidth}\raggedright
accuracy (red)
\end{minipage} & \begin{minipage}[b]{\linewidth}\raggedright
accuracy (blue)
\end{minipage} & \begin{minipage}[b]{\linewidth}\raggedright
accuracy (green)
\end{minipage} \\
\midrule\noalign{}
\endhead
\bottomrule\noalign{}
\endlastfoot
0.0 & 0.895 & 0.928 & 0.925 \\
0.1 & 0.905 & 0.932 & 0.930 \\
0.2 & 0.915 & 0.940 & 0.942 \\
0.3 & 0.925 & 0.945 & 0.948 \\
0.4 & 0.935 & 0.955 & 0.958 \\
0.5 & 0.945 & 0.960 & 0.962 \\
0.6 & 0.955 & 0.970 & 0.972 \\
0.7 & 0.965 & 0.975 & 0.978 \\
0.8 & 0.975 & 0.980 & 0.982 \\
0.9 & 0.985 & 0.985 & 0.988 \\
1.0 & 1.000 & 1.000 & 1.000 \\
\end{longtable}
}

\textit{[Image omitted: images/a0721ec671ccc6654d6f7574e6a996e2fce2438c0599cd666ef242eb351a49ff.jpg]}

line

{\def\LTcaptype{none} % do not increment counter
\begin{longtable}[]{@{}llll@{}}
\toprule\noalign{}
\% of supervised data & MACE-EM & MACE-VB & majority \\
\midrule\noalign{}
\endhead
\bottomrule\noalign{}
\endlastfoot
0.0 & 0.942 & 0.941 & 0.936 \\
0.1 & 0.948 & 0.947 & 0.940 \\
0.2 & 0.952 & 0.951 & 0.946 \\
0.3 & 0.958 & 0.956 & 0.952 \\
0.4 & 0.964 & 0.963 & 0.960 \\
0.5 & 0.972 & 0.971 & 0.969 \\
0.6 & 0.974 & 0.973 & 0.971 \\
0.7 & 0.978 & 0.977 & 0.976 \\
0.8 & 0.984 & 0.983 & 0.982 \\
0.9 & 0.992 & 0.991 & 0.990 \\
1.0 & 1.000 & 1.000 & 1.000 \\
\end{longtable}
}

Figure 7: Varying the amount of supervision: effect on prediction
accuracy. Each point averaged over 10 runs. MACE uses supervision more
efficiently.

they choose to focus on, and the type of annotation task they model. For
example, many methods explicitly model annotator bias in addition to
annotator competence (Dawid and Skene, 1979; Smyth et al., 1995). Our
work models annotator bias, but only when the annotator is suspected to
be spamming.

Other methods focus modeling power on instance difficulty to learn not
only which annotators are good, but which instances are hard (Carpenter,
2008; Whitehill et al., 2009). In machine vision, several models have
taken this further by parameterizing difficulty in terms of complex
features defined on each pairing of annotator and annotation instance
(Welinder et al., 2010; Yan et al., 2010). While such features prove
very useful in vision, they are more difficult to define for the
categorical problems common to NLP. In addition, several methods are
specifically tailored to annotation tasks that involve ranking (Steyvers
et al., 2009; Lee et al., 2011), which limits their applicability in
NLP.

The method of Raykar and Yu (2012) is most similar to ours. Their goal
is to identify and filter out annotators whose annotations are not
correlated with the gold label. They define a function of the learned
parameters that is useful for identifying these spammers, and then use
this function to build a prior. In contrast, we use simple priors, but
incorporate a model parameter that explicitly represents the probability
that an annotator is spamming. Our simple model achieves the same
accuracy on gold label predictions as theirs.

\section{7 Conclusion}\label{conclusion}

We provide a Java-based implementation, MACE, that recovers correct
labels with high accuracy, and reliably identifies trustworthy
annotators. In addition, it provides a threshold to control the
accuracy/coverage trade-off and can be trained with standard EM or
Variational Bayes EM. MACE works fully unsupervised, but can incorporate
token constraints via annotated control items. We show that even small
amounts help improve accuracy.

Our model focuses most of its modeling power on learning trustworthiness
parameters, which are highly correlated with true annotator reliability
(Pearson \(\rho\) 0.9). We show on real-world and synthetic data sets
that our method is more accurate than majority voting, even under
adversarial conditions, and as accurate as more complex state-of-the-art
systems. Focusing on high-confidence instances improves accuracy
considerably. MACE is freely available for download under
http://www.isi.edu/publications/licensed-sw/mace/index.html.

\section{Acknowledgements}\label{acknowledgements}

The authors would like to thank Chris Callison-Burch, Victoria Fossum,
Stephan Gouws, Marc Schulder, Nathan Schneider, and Noah Smith for
invaluable discussions, as well as the reviewers for their constructive
feedback.

\section{References}\label{references}

Jiang Bian, Yandong Liu, Ding Zhou, Eugene Agichtein, and Hongyuan Zha.
2009. Learning to recognize reliable users and content in social media
with coupled mutual reinforcement. In Proceedings of the 18th
international conference on World wide web, pages 51--60. ACM.\\
Chris Callison-Burch and Mark Dredze. 2010. Creating speech and language
data with amazon's mechanical turk. In Proceedings of the NAACL HLT 2010
Workshop on Creating Speech and Language Data with Amazon's Mechanical
Turk, pages 1--12, Los Angeles, June. Association for Computational
Linguistics.\\
Chris Callison-Burch, Philipp Koehn, Christof Monz, Kay Peterson, Mark
Przybocki, and Omar Zaidan. 2010. Findings of the 2010 joint workshop on
statistical machine translation and metrics for machine translation. In
Proceedings of the Joint Fifth Workshop on Statistical Machine
Translation and MetricsMATR, pages 17--53, Uppsala, Sweden, July.
Association for Computational Linguistics.\\
Bob Carpenter. 2008. Multilevel Bayesian models of categorical data
annotation. Unpublished manuscript.\\
A. Philip Dawid and Allan M. Skene. 1979. Maximum likelihood estimation
of observer error-rates using the EM algorithm. Applied Statistics,
pages 20--28.\\
Arthur P. Dempster, Nan M. Laird, and Donald B. Rubin. 1977. Maximum
likelihood from incomplete data via the EM algorithm. Journal of the
Royal Statistical Society. Series B (Methodological), 39(1):1--38.\\
Jason Eisner. 2002. An interactive spreadsheet for teaching the
forward-backward algorithm. In Proceedings of the ACL-02 Workshop on
Effective tools and methodologies for teaching natural language
processing and computational linguistics-Volume 1, pages 10--18.
Association for Computational Linguistics.\\
Alvan R. Feinstein and Domenic V. Cicchetti. 1990. High agreement but
low kappa: I. the problems of two paradoxes. Journal of Clinical
Epidemiology, 43(6):543--549.\\
Kilem Li Gwet. 2008. Computing inter-rater reliability and its variance
in the presence of high agreement. British Journal of Mathematical and
Statistical Psychology, 61(1):29--48.\\
Eduard Hovy. 2010. Annotation. A Tutorial. In 48th Annual Meeting of the
Association for Computational Linguistics.\\
Mukund Jha, Jacob Andreas, Kapil Thadani, Sara Rosenthal, and Kathleen
McKeown. 2010. Corpus creation for new genres: A crowdsourced approach
to pp attachment. In Proceedings of the NAACL HLT 2010 Workshop on
Creating Speech and Language Data with Amazon's Mechanical Turk, pages
13--20. Association for Computational Linguistics.

Mark Johnson. 2007. Why doesn't EM find good HMM POS-taggers. In
Proceedings of the 2007 Joint Conference on Empirical Methods in Natural
Language Processing and Computational Natural Language Learning
(EMNLP-CoNLL), pages 296--305.\\
Michael D. Lee, Mark Steyvers, Mindy de Young, and Brent J. Miller.
2011. A model-based approach to measuring expertise in ranking tasks. In
L. Carlson, C. Hölscher, and T.F. Shipley, editors, Proceedings of the
33rd Annual Conference of the Cognitive Science Society, Austin, TX.
Cognitive Science Society.\\
Vikas C. Raykar and Shipeng Yu. 2012. Eliminating Spammers and Ranking
Annotators for Crowdsourced Labeling Tasks. Journal of Machine Learning
Research, 13:491--518.\\
Padhraic Smyth, Usama Fayyad, Mike Burl, Pietro Perona, and Pierre
Baldi. 1995. Inferring ground truth from subjective labelling of Venus
images. Advances in neural information processing systems, pages
1085--1092.\\
Rion Snow, Brendan O'Connor, Dan Jurafsky, and Andrew Y. Ng. 2008. Cheap
and fast---but is it good? Evaluating non-expert annotations for natural
language tasks. In Proceedings of the Conference on Empirical Methods in
Natural Language Processing, pages 254--263. Association for
Computational Linguistics.\\
Alexander Sorokin and David Forsyth. 2008. Utility data annotation with
Amazon Mechanical Turk. In IEEE Computer Society Conference on Computer
Vision and Pattern Recognition Workshops, CVPRW '08, pages 1--8. IEEE.\\
Mark Steyvers, Michael D. Lee, Brent Miller, and Pernille Hemmer. 2009.
The wisdom of crowds in the recollection of order information. Advances
in neural information processing systems, 23.\\
Stephen Tratz and Eduard Hovy. 2010. A taxonomy, dataset, and classifier
for automatic noun compound interpretation. In Proceedings of the 48th
Annual Meeting of the Association for Computational Linguistics, pages
678--687. Association for Computational Linguistics.\\
Peter Welinder, Steve Branson, Serge Belongie, and Pietro Perona. 2010.
The multidimensional wisdom of crowds. In Neural Information Processing
Systems Conference (NIPS), volume 6.\\
Jacob Whitehill, Paul Ruvolo, Tingfan Wu, Jacob Bergsma, and Javier
Movellan. 2009. Whose vote should count more: Optimal integration of
labels from labelers of unknown expertise. Advances in Neural
Information Processing Systems, 22:2035--2043.\\
Yan Yan, Rómer Rosales, Glenn Fung, Mark Schmidt, Gerardo Hermosillo,
Luca Bogoni, Linda Moy, and

Jennifer Dy. 2010. Modeling annotator expertise: Learning when everybody
knows a bit of something. In International Conference on Artificial
Intelligence and Statistics. Omar F. Zaidan and Chris Callison-Burch.
2011. Crowdsourcing translation: Professional quality from
non-professionals. In Proceedings of the 49th Annual Meeting of the
Association for Computational Linguistics: Human Language Technologies,
pages 1220--1229, Portland, Oregon, USA, June. Association for
Computational Linguistics.

\end{document}
